\subsection{Project Requirements}
\label{sec:requirements}

\begingroup
\fontsize{10pt}{12pt}\selectfont
\ttfamily
\setlength{\parindent}{0pt}
\setlength{\parskip}{6pt}

The verbs ``MUST,'' ``MUST NOT,'' ``SHOULD,'' ``SHOULD NOT,'' and ``MAY'' are to be interpreted per \href{https://datatracker.ietf.org/doc/html/rfc2119}{RFC~2119}.

\textbf{Identity Management and Privacy}

\begin{enumerate}
  \item The system MUST present the option to sign in with Phendo credentials. It MUST also present the option to create a profile for patients who do not use Phendo.
  \item The system MUST inform the patient that all data is stored on-device via a dedicated privacy page.
  \item The system MUST NOT transmit, collect, or share any health data. No server-side storage, analytics, crash reports, or device identifiers SHALL be present.
  \item The system MUST allow the patient to delete all data by clearing local storage or uninstalling the app.
\end{enumerate}

\textbf{Capture}

\begin{enumerate}
  \item The system MUST allow the patient to record their experience through open-ended speech via a dedicated recording interface with a visual waveform. It MUST NOT impose structured input.
  \item The system MUST state that all audio is deleted after transcription.
  \item The system SHOULD use an on-device LLM for speech transcription. (The prototype demonstrates the transcription flow with synthetic data; live on-device transcription is deferred.)
  \item The system MUST extract clinically relevant factors from conversational input using a downstream LLM which MAY be off-device. Extracted factors MUST include:
        \begin{itemize}
          \item Pain (16 locations, 3 severity levels)
          \item Mood (10 positive, 20 negative sentiments)
          \item Period (flow intensity, clots, spotting, breakthrough bleeding)
          \item GI/Urinary symptoms (15 symptoms, 3 severity levels)
          \item Functional impact (``Hard to Do'': 20 activities)
          \item Other symptoms (21 items, 3 severity levels)
          \item Medications (free-text list)
          \item Overall day rating (Good / Manageable / Bad)
        \end{itemize}
  \item Extraction categories MUST align with the Phendo citizen-science taxonomy.
  \item The system MUST surface all extractions transparently as visual badges/pills organized by category so the patient can see what was inferred.
  \item The system MUST allow the patient to override the overall day rating and edit the transcript text. It SHOULD allow the patient to correct or reject individual extracted factors.
  \item The system MUST NOT gamify tracking or create pressure to log consistently.
\end{enumerate}

\textbf{Reflective Sensemaking}

\begin{enumerate}
  \item The system MUST generate per-week and per-month narrative summaries covering pain, mood, period, GI, functional impact, medications, and overall status.
  \item All AI-generated summaries MUST be presented as reflective anchors---observational and patient-centered in tone (e.g., ``Here's what I see in your data'').
  \item The system MUST NOT generate language that implies clinical certainty or medical advice.
  \item Weekly summaries MUST be viewable by month, with collapsible long/short display modes.
  \item Monthly summaries MUST aggregate symptom frequencies and top factors.
  \item The patient SHOULD be able to edit summaries inline.
\end{enumerate}

\textbf{Timeline}

\begin{enumerate}
  \item The system MUST support viewing the illness trajectory at four temporal resolutions: day, week, month, and year.
  \item Day view MUST display a horizontally scrollable date strip organized by month, with each day showing the overall rating and full detail on selection.
  \item Year view MUST render a calendar heatmap colored by category severity.
  \item Year view MUST allow the patient to switch the displayed category among Overall, Pain, Mood, GI, and Period. The selected category SHOULD persist across sessions.
  \item The system SHOULD support arbitrary, user-defined timespans. (Not yet delivered.)
  \item The system SHOULD allow the patient to annotate any point on the timeline with free text and tags. (Not yet delivered.)
\end{enumerate}

\textbf{Visit Preparation}

\begin{enumerate}
  \item The system MUST provide a clinical visit preparation view with interactive anterior and posterior body map visualizations.
  \item Body map zones MUST be clickable and MUST display a summary of mapped symptoms, pain locations, and their frequencies for the selected time range.
  \item Zones MUST be color-coded by maximum severity within the selected time range.
  \item The system MUST support at least four preparation time ranges: last week, two weeks, last month, and six months.
  \item Each time range MUST present a narrative summary covering logging frequency, pain days/locations, period data, GI patterns, functional impact, medications, and mood themes.
  \item The system SHOULD allow the patient to control which data is disclosed to a clinician.
\end{enumerate}

\endgroup
